\documentclass[10pt]{article}

% Standalone layout: compiles in a blank Overleaf project without the
% official ICLR style file. For submission, replace this preamble with
% \usepackage{iclr2027_conference,times} plus the conference template.
\usepackage[letterpaper,margin=1in]{geometry}
\usepackage{times}
\usepackage{microtype}

\usepackage{amsmath,amssymb,amsfonts,amsthm,mathtools,bm}
\usepackage{enumitem}
\usepackage{booktabs}
\usepackage{xcolor}
\usepackage{url}
\usepackage{algorithm}
\usepackage{algpseudocode}

\newtheorem{assumption}{Assumption}
\newtheorem{theorem}{Theorem}
\newtheorem{lemma}{Lemma}
\newtheorem{remark}{Remark}
\newtheorem{corollary}{Corollary}
\newtheorem{definition}{Definition}

\newcommand{\E}{\mathbb{E}}
\newcommand{\X}{\mathcal{X}}
\newcommand{\F}{\mathcal{F}}
\newcommand{\argmin}{\mathop{\mathrm{argmin}}}
\newcommand{\dist}{\mathop{\mathrm{dist}}}

\title{Proximal Drift-Plus-Penalty for Weakly Convex Optimization\\
with Time-Average Constraints: A Soft Prompting Motivation}
\author{Anonymous Authors}

\begin{document}
\maketitle


\begin{abstract}
Soft prompting provides a concrete setting in which an optimization
objective must be balanced against operational constraints that are
naturally imposed over a long horizon. We consider a stationary
prompt-selection policy and use this setting to motivate weakly convex
optimization with one ordinary constraint and one long-run constraint that
may itself change from round to round. We develop a proximal
drift-plus-penalty (DPP) method in which the long-run constraint is
stabilized before queue weighting. The stabilization makes the
queue-weighted constraint component strongly convex and yields a proximal
subproblem with a queue-independent baseline curvature. We prove a
pointwise $O(V)$ virtual-queue bound, an $O(V/T)$ time-average constraint
guarantee, and an $O(1/V+1/T+V_T/T)$ bound on the squared
near-stationarity residual of the penalized problem, where $V_T$ is the
cumulative functional variation of the per-round constraint loss. With
additional Lipschitz and violating-point regularity assumptions, we
convert this penalized stationarity guarantee into an
$O(1/V+1/T+V_T/T)$ expected violation of the ordinary constraint for a
uniformly sampled iterate. The resulting statement is deliberately a
nearly-KKT guarantee rather than a claim of convergence of the entire
nonconvex iterate sequence to an exact KKT point.
\end{abstract}

\section{Problem Formulation and Assumptions}
\label{sec:problem}

We study the constrained problem
\begin{equation}
\min_{x\in\X} f(x)
\quad\text{s.t.}\quad
g(x)\leq0,
\qquad
\limsup_{T\to\infty}\frac1T\sum_{t=1}^T h_t(x_t)\leq0,
\label{eq:main-problem}
\end{equation}
where $h_t$ is the constraint loss in force at round $t$: if $x_t$ denotes
the decision used at round $t$, the operational requirement is on the
\emph{time average} of the sequence $h_1,h_2,\ldots$ evaluated along the
algorithm's own trajectory. Nothing below assumes $h_t$ is the same
function from one round to the next.

We enforce the ordinary constraint through the exact hinge penalty
\begin{equation}
P(x)=f(x)+\beta[g(x)]_+,
\qquad
[a]_+:=\max\{a,0\}.
\label{eq:P}
\end{equation}

\begin{assumption}[Decision set and boundedness]
\label{ass:compact}
$\X$ is nonempty, closed, convex, and compact, with diameter
$\sup_{x,y\in\X}\|x-y\|\leq D$. The functions $f,g$, and every $h_t$ are
finite on an open neighborhood of $\X$.
\end{assumption}

\begin{assumption}[Weak convexity]
\label{ass:weak}
$f$ is $\rho_f$-weakly convex, $g$ is $\rho_g$-weakly convex, and for
every $t$, $h_t$ is $\rho_h$-weakly convex on a neighborhood of $\X$.
Since $[g]_+$ is also $\rho_g$-weakly convex, $\rho_P:=\rho_f+\beta\rho_g$
is a valid weak-convexity modulus for $P$.
\end{assumption}

\begin{assumption}[Lipschitz continuity]
\label{ass:lipschitz}
$f,g$, and every $h_t$ are respectively $L_f,L_g,L_h$-Lipschitz on $\X$.
\end{assumption}

\begin{assumption}[Bounded constraint values]
\label{ass:boundedh}
For all $t$ and all $x\in\X$, $|h_t(x)|\leq H$.
\end{assumption}

\begin{assumption}[Strict time-average feasibility]
\label{ass:slater}
There exists a point $x^{\circ}\in\X$ and $\xi>0$ such that
$h_t(x^{\circ})\leq-\xi$ for all $t$.
\end{assumption}

This is a Slater-type condition for the time-average constraint, required
to hold at a single common comparison point across every round.
\textbf{Importantly, $x^{\circ}$ need not minimize $P$} --- it is used
only as a comparison point in the queue argument, and the proof of
Lemma~\ref{lem:queuepoint} below drops $P(x_{t+1})$ only as far as its
global lower bound $P_{\min}$, never as far as $P(x^{\circ})$ itself. The
strict margin is needed because the stabilization term
$\gamma\|x-x_t\|^2$ consumes part of it.

\begin{assumption}[Violating-point regularity for the ordinary constraint]
\label{ass:vp}
There exists $\kappa_g>0$ such that, for every $x\in\X$ satisfying
$g(x)>0$,
\[
\dist\bigl(0,\ \beta\,\partial g(x)+N_{\X}(x)\bigr)\geq\kappa_g.
\]
For a smooth constraint and an interior point of $\X$, this reduces to
$\beta\|\nabla g(x)\|\geq\kappa_g$ whenever $g(x)>0$.
\end{assumption}

\begin{assumption}[Algorithmic parameters]
\label{ass:params}
Choose
\[
\frac{\rho_h}{2}<\gamma<\frac{\xi}{D^2},
\qquad
0<\eta<\frac1{\rho_P}
\]
(any $\eta>0$ when $\rho_P=0$), and let $V\geq1$. Initialize $Q_1=0$.
\end{assumption}

\section{Algorithm}
\label{sec:algorithm}

At round $t$, only $h_t$ --- the constraint function \emph{currently in
force} --- is observed and used; nothing about $h_{t+1}, h_{t+2},\ldots$
is required in advance.

\begin{algorithm}[h]
\caption{Proximal Drift-Plus-Penalty}
\label{alg:pdpp}
\begin{algorithmic}[1]
\Require Initial point $x_1\in\X$, queue $Q_1=0$, parameters $V\geq1$,
$\eta>0$, $\gamma>0$
\For{$t=1,\ldots,T$}
    \State Observe the current round's constraint function $h_t$
    \State Compute
    \[
    x_{t+1}\in\argmin_{x\in\X}
    \left\{
    P(x)+\frac{Q_t}{V}\left[h_t(x)+\gamma\|x-x_t\|^2\right]
    +\frac1{2\eta}\|x-x_t\|^2
    \right\}
    \]
    \State Update the virtual queue
    \[
    Q_{t+1}\gets[Q_t+h_t(x_{t+1})]_+
    \]
\EndFor
\end{algorithmic}
\end{algorithm}

\begin{lemma}[Uniform well-posedness]
\label{lem:wellposed}
Under Assumptions~\ref{ass:compact}--\ref{ass:params}, the subproblem in
Algorithm~\ref{alg:pdpp} is strongly convex with modulus
\[
\mu_t=\frac1\eta-\rho_P+q_t(2\gamma-\rho_h)\geq\mu_0:=\frac1\eta-\rho_P>0,
\qquad q_t:=\frac{Q_t}{V},
\]
for every $t$, regardless of which $h_t$ is in force that round.
Consequently $x_{t+1}$ is unique.
\end{lemma}

\section{Pointwise Queue Bound}
\label{sec:guarantees}

Let $P_{\min}:=\min_{x\in\X}P(x)$, $P_{\max}:=\max_{x\in\X}P(x)$, and
define
\[
C_\circ:=P(x^{\circ})-P_{\min}+\frac{D^2}{2\eta},
\qquad
\bar\xi:=\xi-\gamma D^2>0,
\qquad
\bar q:=\frac{2C_\circ}{\bar\xi}+H.
\]

\begin{lemma}[Pointwise queue bound]
\label{lem:queuepoint}
Under Assumptions~\ref{ass:compact}--\ref{ass:params},
$Q_t\leq 2VC_\circ/\bar\xi+H$ for all $t$. Consequently, for $V\geq1$,
$q_t\leq\bar q$ for all $t$.
\end{lemma}

This holds already under Assumptions~\ref{ass:compact}--\ref{ass:params}
as stated --- each $h_t$ may differ from round to round, so long as it
satisfies the \emph{uniform} moduli/bound/margin those assumptions
require of \emph{every} $t$.

\section{Functional Variation and the Movement/Stationarity Bounds}
\label{sec:tv}

The queue bound above does not depend on how much $h_t$ changes from one
round to the next. The movement bound does: a lower bound on how
$h_{t+1}$ differs from $h_t$ at the same point $x_{t+1}$ appears when the
per-round potential is telescoped, and it must be paid for explicitly.

\begin{definition}[Functional variation]
\label{def:variation}
Define the one-step functional variation
$\Delta_t:=\sup_{x\in\X}|h_{t+1}(x)-h_t(x)|$
and the cumulative variation $V_T:=\sum_{t=1}^{T-1}\Delta_t$; write
$\nu_T:=V_T/T$ for the average functional variation.
\end{definition}

\begin{assumption}[Uniformity across rounds]
\label{ass:tv}
The moduli $\rho_h,L_h,H$ in Assumptions~\ref{ass:weak}--\ref{ass:boundedh}
and the point $x^{\circ}$ and margin $\xi$ in
Assumption~\ref{ass:slater} are common to every $t$ (not merely existing
separately for each $t$).
\end{assumption}

\begin{lemma}[Well-posedness and pointwise queue bound, restated]
\label{lem:tv-queue}
Under Assumption~\ref{ass:tv}, Lemma~\ref{lem:wellposed} and
Lemma~\ref{lem:queuepoint} hold with the same constants $\mu_0$,
$C_\circ$, $\bar\xi$, $\bar q$ at every $t$, uniformly over the whole run.
\end{lemma}

\begin{theorem}[Long-run constraint feasibility]
\label{thm:feasibility}
Under the assumptions above,
\begin{equation}
\frac1T\sum_{t=1}^T h_t(x_{t+1})
\leq
\frac{2VC_\circ}{\bar\xi T}+\frac HT.
\label{eq:avg-feas}
\end{equation}
Consequently, the time-average constraint violation is $O(V/T)$.
\end{theorem}

\begin{theorem}[Movement and nearly-stationary guarantee]
\label{thm:movement}
Define $c_0:=1/\eta-\rho_P/2>0$, $C_S:=P_{\max}-P_{\min}+2\bar qH$,
$A:=1/\eta+2\gamma\bar q$. Let $\tau$ be uniform on $\{1,\ldots,T\}$, and
define
\[
R_t:=\dist\bigl(0,\ \partial P(x_{t+1})+q_t\,\partial h_t(x_{t+1})
+N_{\X}(x_{t+1})\bigr).
\]
Then
\begin{equation}
\frac1T\sum_{t=1}^T\|x_{t+1}-x_t\|^2
\leq
\frac{C_S}{c_0T}+\frac{H^2}{c_0V}+\frac{\bar q\,V_T}{c_0T}+\frac{D^2}{T},
\label{eq:movement-tv}
\end{equation}
and
\begin{equation}
\E[R_\tau^2]\leq A^2\cdot(\text{the right-hand side of \eqref{eq:movement-tv}})
=O\!\left(\frac1V+\frac1T+\frac{V_T}{T}\right).
\label{eq:residual-tv}
\end{equation}
Consequently $\E[R_\tau]=O\bigl(V^{-1/2}+T^{-1/2}+\sqrt{V_T/T}\bigr)$. If
$h_t\equiv h$ is in fact constant in $t$, then $V_T=0$ and this recovers
the time-invariant movement bound exactly.
\end{theorem}

\begin{theorem}[Nearly-feasible stationarity for the ordinary constraint]
\label{thm:gfeas}
Suppose in addition that
$\dist(0,\beta\partial g(x)+N_{\X}(x))\geq\kappa_g$ whenever $g(x)>0$,
that $L_f,L_h$ are uniform subgradient bounds, and define
$\delta_g:=\kappa_g-L_f-\bar qL_h>0$,
$G_g:=\sup_{x\in\X}[g(x)]_+$. Then
\begin{equation}
\E[g(x_{\tau+1})]_+
\leq
\frac{G_g}{\delta_g^2}\,\E[R_\tau^2]
=
O\!\left(\frac1V+\frac1T+\frac{V_T}{T}\right).
\label{eq:gfeas-tv}
\end{equation}
\end{theorem}

\begin{remark}[Interpretation]
The temporal variation contributes only through $V_T/T$ in the squared
stationarity residual. Thus $V_T=o(T)$ is sufficient for the variation
penalty to vanish, while $V_T=O(1)$ recovers the stationary-order bound.
A constraint that changes by a bounded total amount over the whole
horizon costs nothing asymptotically; one that keeps drifting at a
constant per-round rate ($V_T=\Theta(T)$) leaves a floor in the movement
and residual bounds that $T\to\infty$ alone cannot remove.
\end{remark}

\bigskip
\noindent\textit{Note: proofs of Lemma~\ref{lem:wellposed},
Lemma~\ref{lem:queuepoint}, and
Theorems~\ref{thm:feasibility}--\ref{thm:gfeas} were not part of the
pasted text this draft was reconstructed from, so none are included here.
A structurally equivalent time-invariant version of this paper (shown
earlier in the same conversation) does contain full appendix proofs,
built on the Lyapunov function
$W_t=\tfrac12Q_t^2+\tfrac{V}{2\eta}\|x_t-x^{\circ}\|^2$; adapting them to
$h_t$ only requires tracking the extra term
$q_{t+1}(h_{t+1}(x_{t+1})-h_t(x_{t+1}))\le\bar q\,\Delta_t$ that appears
when the per-round potential is telescoped across a round where the
constraint function itself changed.}

\end{document}
